\documentclass{jsarticle}
\usepackage[dvipdfmx]{graphicx}
\usepackage{chemfig}
\usepackage{amsmath}
\usepackage{amssymb}
\begin{document}
\title{反重力発生装置の各構成物質}
\author{Masaaki Yamaguchi}
\date{}
\maketitle



  原子振動子　セシウム Cs


  形状記憶合金 (Fe・Co60・Pt・Al)$H_2 SO_4, Al_2 O_3$


  反重力発生器 Pr(パラジウム合金)電磁場生成He,$H_2$Mg,Al(摩擦熱)


  結晶石 Co60,Pt,Prの反陽子


  緩衝剤 慣性力感知器 有機化合物とS化合物(シクロアルカン$C_n H_{2n}$
  方位と位置探知機にも使う


  量子コンピュータの量子素子 Pt,Ag,Au（電子と電流のエネルギー経路)半導体に
  使う
  回路の合成演算子 


  $$ x^y = {1 \over {y^x}} \to $$

  $$ \pi(\chi,x) = [i\pi(\chi,x),f(x)] $$
  $$ {d \over df}F = F^{{f}^{'}} $$


  ディスプレイの電磁迷彩 (Al,Mg) $\to$ S合成子 $Si O_2 $
  プリズム C

\end{document}
